\section{Summary of analytical solution}\label{mathystuff}

To fix ideas, consider the day a bud blooms at a given location. We
model the effective temperature experienced by that bud on day \(i\) by
a linear trend plus microclimate noise, \begin{equation}
X_i = \alpha + \beta\, i + \epsilon_i,
\end{equation} where \(\alpha\) is a baseline level, \(\beta\) is a
within-season warming rate, and \(\{\epsilon_i\}_{i\ge 1}\) are
independent and identically distributed with
\(\mathbb{E}[\epsilon_i]=0\) and \(\mathrm{Var}(\epsilon_i)=\sigma^2\),
representing idiosyncratic variation at that bud such as due to the
day-to-day microclimate.

Let \begin{equation}
Z_n \;=\; \sum_{i=1}^n X_i
\qquad\text{and}\qquad
\nu(t) \;=\; \min\{m:\, Z_m>t\}
\end{equation} denote cumulative temperature by day \(n\) and the
observed bloom date, respectively, where \(t>0\) is a plant-specific
tempearture-sum threshold. Thus \(\nu(t)\) is the hitting time at which
cumulative temperatures first exceed the threshold.

Standard results for stopped random walks yield asymptotic normal
approximations for \(\nu(t)\), with behavior governed by the parameters
\(\alpha\) and \(\beta\). We consider two cases.

\paragraph{Regime 1: constant drift ($\beta=0$)}

When the drift is approximately constant and positive, \begin{equation}
\nu(t)\ \dot{\sim}\ \mathrm{Normal}\!\left(\frac{t}{\alpha},\ \frac{\sigma^2\,t}{\alpha^{3}}\right),
\label{eq:nu_beta0}
\end{equation} see Gut (2011) and Gut (2009, chap. 3).

\paragraph{Regime 2: increasing drift ($\beta>0$)}

When average temperatures increase, the hitting time satisfies
\begin{equation}
\nu(t)\ \dot{\sim}\ \mathrm{Normal}\!\left(\sqrt{\frac{2t}{\beta}} - \frac{\alpha}{\beta} + \frac{1}{2},\ \frac{\sigma^2}{\beta^{3/2}\sqrt{2t}}\right),
\label{eq:nu_beta_pos}
\end{equation} with the derivation sketched in Appendix A.1 and
simulations confirming the derivations in Appendix A.3. The functional
central limit theorem and linearization argument underlying these
results hold somewhat generally, and thus similar limits can be derived
when the \(\{\epsilon_i\}_{i\ge 1}\) are weakly correlated or the
temperature trajectory is not linear but sufficiently.

From equations \eqref{eq:nu_beta0}--\eqref{eq:nu_beta_pos}, we make two
observations. Finding 1: there is a nonlinear relationship between
warming and the mean timing of springtime events, and warming can
advance springtime events at an increasing rate. Finding 2: the two
regimes are not equally sensitive to the idiosyncratic variation, and
thus warming can increase the variation of springtime events, thereby
scattering spring. We now discuss both findings in greater detail.

\subsection{2.1 Finding 1: warming can advance springtime events at an
increasing
rate}\label{finding-1-warming-can-advance-springtime-events-at-an-increasing-rate}

In both regimes, the expected bloom date is a nonlinear function of the
warming parameters.

\paragraph{Regime 1: constant drift ($\beta=0$)}

From \eqref{eq:nu_beta0}, \begin{equation}
\mathbb{E}[\nu(t)] \approx \frac{t}{\alpha},
\qquad
\mathrm{Var}(\nu(t)) \approx \frac{\sigma^2\,t}{\alpha^3}.
\end{equation} Increasing baseline temperatures (larger \(\alpha\))
advances the expected bloom date, but the marginal advance diminishes
since \begin{equation}
\frac{\partial}{\partial \alpha}\,\mathbb{E}[\nu(t)] \approx -\frac{t}{\alpha^2},
\end{equation} where the magnitude decreases as \(\alpha\) grows.

\paragraph{Regime 2: increasing drift ($\beta>0$)}

From \eqref{eq:nu_beta_pos}, when \(t\) is large, \begin{equation}
\mathbb{E}[\nu(t)] \approx \sqrt{\frac{2t}{\beta}},
\qquad
\mathrm{Var}(\nu(t)) \approx \frac{\sigma^2}{\beta^{3/2}\sqrt{2t}}.
\end{equation} Increasing the within-regime warming rate (larger
\(\beta\)) also advances the expected bloom date with diminishing
returns since \begin{equation}
\frac{\partial}{\partial \beta}\,\mathbb{E}[\nu(t)]
\approx
-\frac{1}{2}\sqrt{\frac{2t}{\beta^{3}}}.
\end{equation} However, if warming climates shift the bloom date from
regime 2 to regime 1, then the rate of advancement can increase,
particularly if increases in \(\alpha\) outpace increases in \(\beta\).

\subsection{2.2 Finding 2: warming can scatter spring
events}\label{finding-2-warming-can-scatter-spring-events}

Equations \eqref{eq:nu_beta0}--\eqref{eq:nu_beta_pos} also clarify how
warming affects dispersion. Holding the regime fixed, warming reduces
variance. That is, \(\mathrm{Var}(\nu(t))\) decreases with \(\alpha\) in
\eqref{eq:nu_beta0} and with \(\beta\) in \eqref{eq:nu_beta_pos}.
However, just as in Section 2.1, warming can change which regime is
relevant at the realized bloom date. When an event is advanced
sufficiently early, microclimate fluctuations play a larger role in
determining the hitting time.

One way to quantify sensitivity to microclimate variation is through the
(squared) coefficient of variation, \[
\mathrm{CV}^2 \;=\; \frac{\mathrm{Var}(\nu(t))}{\mathbb{E}[\nu(t)]^2}.
\] In regime 1, \(\mathbb{E}[\nu(t)]\approx t/\alpha\) and
\(\mathrm{Var}(\nu(t))\approx \sigma^2 t/\alpha^3\), so
\begin{equation*}
\mathrm{CV}^2_{1}
\;=\;
\frac{\mathrm{Var}(\nu(t))}{\mathbb{E}[\nu(t)]^2}
\;\approx\;
\frac{\sigma^2 t/\alpha^3}{(t/\alpha)^2}
\;=\;
\frac{\sigma^2}{\alpha\,t}.
\end{equation*} In regime 2,
\(\mathbb{E}[\nu(t)]\approx \sqrt{2t/\beta}\) and
\(\mathrm{Var}(\nu(t))\approx \sigma^2/(\beta^{3/2}\sqrt{2t})\), hence
\begin{equation*}
\mathrm{CV}^2_{2}
\;=\;
\frac{\mathrm{Var}(\nu(t))}{\mathbb{E}[\nu(t)]^2}
\;\approx\;
\frac{\sigma^2/(\beta^{3/2}\sqrt{2t})}{(2t/\beta)}
\;=\;
\frac{\sigma^2}{\sqrt{\beta} \, (2t)^{3/2}}.
\end{equation*} Both expressions increase when deterministic
accumulation is weaker relative to the microclimate scale \(\sigma\)
(e.g., smaller \(\alpha\) and/or smaller effective \(\beta\)).
Consequently, warming tends to reduce variability within a given
seasonal regime by strengthening deterministic forcing.

To compare regimes, note that \begin{equation*}
\frac{\mathrm{CV}^2_{1}}{\mathrm{CV}^2_{2}}
\;\approx\;
\frac{\sigma^2/(\alpha t)}{\sigma^2/(\sqrt{\beta} \, (2t)^{3/2})}
\;=\;
\frac{\sqrt{\beta} \, 2^{3/2}}{\alpha}\,\sqrt{t}.
\end{equation*} Thus, since \(\alpha\) and \(\beta\) are nearly always
greater than \(0\), the regime 1 becomes arbitrarily more sensitive to
microclimate fluctuations than regime 1 for large \(t\). This produces a
more variable, less synchronized and thus ``scattered'' spring when
warming shifts threshold crossings from regime 2 to regime 1.

\subsection{Experimental Data}\label{experimental-data}

We first consider a controlled experiment on walnut trees
(\textit{Juglans}\textasciitilde sp.) conducted by Charrier et al.
(2011). The data examines 30 trees comprising 6 genotypes at 2
locations. In November of the study year, stems were sampled from each
tree and cut into 7 centimeter segments containing a single bud. All
segments were chilled at \(4^\circ\)C and then transferred to constant
``forcing'' environments with temperatures set at one of \(5,10,15,20,\)
or \(25^\circ\)C. For each temperature, the outcome is the response
time, the time (in days) until budburst determined from the date at
which 50\% of buds reached stage 15 of the BBCH scale.

This experiment isolates the winter regime because temperature is held
constant over time. Recall that in such cases, cumulative forcing grows
approximately linearly, \[
Z_n \approx \alpha n + \sum_{i=1}^n e_i,
\] so that the hitting time \(\nu(t)\) satisfies the approximation \[
\nu(t)\ \dot{\sim}\ \mathrm{Normal}\!\left(\frac{t}{\alpha},\ \frac{\sigma^2 t}{\alpha^3}\right),
\] implying that increasing \(\alpha\) should (i) reduce the mean time
to budburst and (ii) reduce dispersion. In addition, the squared
coefficient of variation, \[
\mathrm{CV}^2=\frac{\mathrm{Var}(\nu(t))}{\mathbb{E}[\nu(t)]^2},
\] provides a unit-free measure of sensitivity to stochastic
fluctuations, which should decrease as \(\alpha\) increases.

Table \ref{tab:walnut_forcing} reports the mean and standard deviation
of the response times across the 30 stems at each forcing temperature,
along with the sensitivity measure \(\mathrm{CV}^2\). The pattern is
consistent with the regime 1 model. Higher forcing temperatures
substantially advance budburst (the mean falls sharply with \(\alpha\)),
and the dispersion of budburst timing is smaller at warmer forcing
temperatures. Thus, in an experimental setting, the stopped-random-walk
approximation captures the basic comparative statics of
constant-temperature forcing on both the level and variability of
budburst times. (The sensitivty does decline from \(\alpha = 20\) to
\(\alpha = 25\). The reason for low sensitivity for \(\alpha = 5\) and
\(10\) is unclear and may reflect censoring.)

\begin{table}[t]
\centering
\caption{Summary data from the Walnut forcing experiment [@charrier2011budburst].}
\label{tab:walnut_forcing}
\begin{tabular}{rrrrr}
\toprule
$\alpha$ & $n$ & mean & sd & $\mathrm{cv}^2$ \\
\midrule
 5  & 30 & 178 & 27.28 & 0.0236 \\
10  & 30 &  88 & 12.44 & 0.0199 \\
15  & 30 &  39 & 11.33 & 0.0831 \\
20  & 30 &  26 &  7.67 & 0.0851 \\
25  & 30 &  20 &  4.45 & 0.0480 \\
\bottomrule
\end{tabular}
\end{table}